\section{Related Work}
\label{sec:related_work}

\subsection{Spherical vs.\ Flat Embedding Spaces}
\label{sec:sphere_vs_flat}

SSL methods largely split into two camps by where they place their learned features. \emph{Spherical} methods---SimCLR, DINO/DINOv2, BYOL, MoCo---apply an $\ell_2$ normalization at the end of the projection head, constraining features to the unit hypersphere. \emph{Flat-space} methods---VICReg, Barlow Twins, SimDINO---leave features unnormalized in Euclidean space and prevent scale drift via distribution regularization. Each choice carries genuine tradeoffs, and understanding them is necessary for understanding why a flat-space objective needs to be constructed the way we construct ours.

\paragraph{What spherical methods gain.}
The hypersphere is a homogeneous space with a canonical metric: angular distance, or equivalently cosine similarity. Because every point has unit norm, ``direction'' is the only remaining degree of freedom, and this radical simplification has real consequences. A linear classifier on the sphere reduces to a hyperplane through the origin, which intersects the sphere in a great circle: the decision boundary is unambiguous, interpretable as a direction, and calibrated in a meaningful metric. This is the geometric foundation of angular-margin face recognition \citep{wang2017normface, meng2021magface, kim2022adaface}, of prototype-based SSL methods where class identities are direction vectors \citep{caron2021emerging, oquab2023dinov2}, and of the contrastive-loss family where cosine similarity against negatives is the natural objective \citep{chen2020simclr}. $\ell_2$ normalization also gives training stability for free: features cannot drift to arbitrarily large or small magnitudes, so temperature scaling, softmax, and other bounded nonlinearities are well-posed without additional scale anchors.

\paragraph{What flat methods gain.}
The cost of sphericalization is that a degree of freedom---the feature norm---is thrown away. Recent work on feature norms in non-normalized SSL models shows that this discarded magnitude is not noise: it carries a model-confidence or typicality signal. Samples the encoder finds easier to recognize end up with larger norms than atypical or corrupted samples, without any supervision to that effect \citep{draganov2025embedding_norms, kacker2025normconfidence}. Analogous observations predate SSL: in word embeddings, vector magnitude correlates with term specificity and frequency, and in face recognition with unnormalized features, magnitude correlates with image quality \citep{meng2021magface}. Flat-space SSL is also the natural setting for operations that combine features semantically---addition, subtraction, interpolation, linear projection---where the magnitude of a result is a meaningful feature of the operation rather than something to be normalized away.

\paragraph{The price of flatness.}
The extra degree of freedom in flat space must be paid for. The encoder's output can in principle drift to arbitrarily large or small magnitudes, collapse to a constant, or concentrate along a single direction. These failure modes have to be blocked explicitly, which is why every flat-space SSL method includes one or more regularizers targeting the feature distribution's second-order moments: variance along each coordinate, correlations between coordinates, or the full covariance matrix. The choice of regularizer is precisely where existing methods diverge from one another, and this paper's central contribution is to argue that the diverse-looking choices in the literature are approximations of a common underlying object.

\subsection{Flat-Space Self-Supervised Methods}

Three methods dominate the flat-space SSL literature. We describe each qualitatively; precise forms and the relationships between them appear in \secref{sec:method}.

\paragraph{VICReg} \citep{bardes2022vicreg} trains a plain Siamese network with three loss terms applied on the projector output. An \emph{invariance} term minimizes the squared-Euclidean distance between features of two augmented views. A \emph{variance} term applies an asymmetric hinge to the per-dimension standard deviation, penalizing directions whose variance falls below one but letting directions above one pass without cost. A \emph{covariance} term penalizes the Frobenius norm of the off-diagonal covariance matrix, encouraging decorrelation. The asymmetric hinge is a deliberate design choice: informative directions that want more variance are not suppressed, protecting the representation from over-flattening.

\paragraph{Barlow Twins} \citep{zbontar2021barlow} takes a different structural approach. Instead of regularizing within-view covariance, it works with the cross-view cross-correlation matrix: each view is batch-normalized, and the inner product between the two normalized batches gives a matrix whose diagonal encodes per-dimension cross-view agreement (invariance plus scale control) and whose off-diagonal encodes inter-dimension redundancy (decorrelation). A single loss drives the diagonal to one and the off-diagonal to zero, fusing what VICReg separates into three terms. Because Barlow's loss operates on batch-normalized features, the per-coordinate scale is implicitly pinned---scale and decorrelation are not free parameters to be balanced, they enter the loss together.

\paragraph{SimDINO} \citep{wu2025simdino} starts from DINO and asks what happens if the prototype-softmax head is replaced with the simplest possible alignment objective. The resulting recipe is a squared-Euclidean alignment loss to an EMA teacher, paired with a \emph{coding-rate} term that encourages the feature distribution to spread out in the sense of maximum-coding-rate-reduction (MCR\textsuperscript{2}) \citep{yu2020mcr2, ma2007segmentation}. The coding-rate term is motivated by a rate-distortion argument: intuitively, it approximates the bits required to describe the batch feature matrix up to a specified distortion. As originally proposed, \textbf{SimDINO is a spherical method}: features are $\ell_2$-normalized to the unit hypersphere before the loss is computed, inheriting the bounded-feature-space stability that every other spherical SSL method enjoys. The coding-rate term on hyperspherical features is well-behaved because the feature magnitudes are fixed at one.

SimDINO enters our paper because the coding-rate term is the closest existing object to what we propose, and because \emph{unnormalized} variants of SimDINO-style training have appeared in the recent flat-space SSL literature. These variants face the scale-stabilization problem discussed in \secref{sec:coding_rate_scale}: the coding-rate term on raw unnormalized covariance does not have an intrinsic scale, and training diverges without a separate scale anchor. A workaround---switching from \emph{coding rate on covariance} to \emph{coding rate on the correlation matrix}---has been proposed \citep{kacker2025correlation_rate}, making the rate term scale-invariant. But scale-invariance resolves only half of the problem: coding rate on covariance grows unboundedly as features are scaled up (wanting infinite spread), while coding rate on correlation is indifferent to scale entirely (having no preferred magnitude), and neither one pins the scale to a specific target. Unnormalized SimDINO variants therefore pair their rate term with an explicit scale regularizer, typically SigReg \citep{schaeffer2024sigreg}, a two-sided quadratic penalty on per-dimension standard deviation. The split-recipe structure---rate plus scale---is a direct consequence of the expansion/matching distinction we develop in \secref{sec:code_rate_expansion}.

\paragraph{The common structural recipe.}
Although they look quite different at first glance---asymmetric hinge vs.\ symmetric quadratic, Frobenius covariance vs.\ coding rate, within-view vs.\ cross-view---these methods agree on a three-part structure: an invariance term linking the two views, a scale-control mechanism keeping per-dimension variance finite and bounded away from zero, and a decorrelation mechanism suppressing off-diagonal structure. Each method chooses its own functional form for each ingredient and tunes the weights between them empirically. We take this coincidence of structure to be evidence that the methods are all targeting the same underlying geometric object, and the diversity of concrete formulas to be evidence that no one has identified the object cleanly. Identifying it is the subject of \secref{sec:method}.

\subsection{Coding Rate and the Expansion-vs-Matching Distinction}
\label{sec:code_rate_expansion}

The coding-rate framework of MCR\textsuperscript{2} \citep{yu2020mcr2} is worth treating separately because it is the closest existing literature to our objective in form, and differs from ours in a structurally important way.

MCR\textsuperscript{2} defines a rate-distortion-style quantity that measures how many bits are needed to encode a batch of features, and argues that representation learning should \emph{maximize} this quantity (the representation should be maximally informative). Operationally, this takes the form of a log-determinant term on the batch covariance that grows with the spread of the features. SimDINO and SimDINOv2 adopt this term as their non-collapse regularizer.

The key distinction we want to draw is between \emph{expansion} and \emph{matching} objectives. An expansion objective rewards features for being more spread out; formally, it is monotone increasing in the spectrum of the covariance and has no upper bound. Coding rate is an expansion objective: there is no finite optimum; larger features are always better by this criterion. A matching objective, by contrast, has a target distribution in mind and is zero only when the empirical distribution agrees with the target. Objectives that match features to a specific second-order target---unit variance per dimension, zero correlation---are matching objectives: they have true minima, not just gradient pressure in a preferred direction.

This distinction has practical consequences. An expansion objective, used alone, cannot stably determine feature scale---it rewards scaling features up indefinitely. Training stability requires either explicit scale control (by $\ell_2$ normalization or by a separate regularizer like SigReg) or an implicit scale-bounding term like squared-Euclidean alignment, which grows with feature magnitude. In practice the scale that emerges is set by whichever term happens to win the scale-balance tug-of-war, which depends on the ratio of loss weights and is therefore sensitive to hyperparameter choices. A matching objective avoids this by building the scale anchor into itself. We develop this argument quantitatively in \secref{sec:coding_rate_scale}.

\subsection{Whitening and Second-Order Matching}

Whitening---the operation of transforming a batch of features so that their empirical covariance is the identity---has been studied in several SSL variants. W-MSE \citep{ermolov2021wmse} applies a closed-form whitening transformation inside the loss, and subsequent work has examined the effects of inserting batch-whitening layers between the backbone and the projector \citep{hua2021shuffled}. The objective our paper proposes shares whitening's target---covariance equal to identity---but differs in mechanism. Whitening is an operation: after it runs, features are whitened exactly, with no trade-off against other objectives. Our objective is a gradient signal: it pushes features toward whitening but lets the optimizer balance that pressure against invariance and against whatever representational structure emerges naturally. These two design choices lead to different spectra at convergence: hard whitening produces exactly unit-variance features, while soft whitening as a loss term produces approximately-unit-variance features with some residual spread.

A parallel literature studies the spectrum of SSL-trained features as a diagnostic for representation quality. RankMe \citep{garrido2023rankme} uses the effective rank of the feature matrix as a predictor of downstream linear-probe performance; $\alpha$-ReQ \citep{agrawal2022alphareq} characterizes the eigenspectrum decay rate as a similar proxy. These are diagnostic: they measure what happens to the spectrum, without prescribing what it should be. Our objective is prescriptive: it specifies a target spectrum ($\lambda_i = 1$ for all $i$) and lets the training dynamics pursue that target.

\subsection{Our Position}

The three flat-space methods described above all regularize the second-order structure of the batch feature distribution, each with its own formulas and its own hyperparameter-balance decisions. The whitening literature targets the same second-order structure via a hard transformation rather than a soft loss. The coding-rate framework provides an information-theoretic motivation for expansion-style regularizers but requires pairing with separate scale control. We argue that all four threads---scale pinning, decorrelation, whitening, and coding rate---are converging on the same object: a divergence between the empirical feature covariance and the identity covariance. The existing recipes are all approximations to or variations of that divergence. Identifying the unified object explicitly, deriving the existing recipes as special cases or approximations, and training directly on the unified object is the contribution of the remainder of this paper.
